\documentclass[11pt]{article}

\usepackage[margin=1in]{geometry}
\usepackage{amsmath,amssymb}
\usepackage{booktabs}
\usepackage{hyperref}
\usepackage[numbers,sort&compress]{natbib}

\title{Mean Splitting is Hard to Beat for Univariate Distributional Forests}
\author{Silas Koemen}
\date{\today}

\newcommand{\drf}{\mathrm{DRF}}
\newcommand{\mmd}{\mathrm{MMD}}
\newcommand{\cart}{\mathrm{CART}}

\begin{document}
\maketitle

\begin{abstract}
Distributional random forests replace mean-based CART splitting with
distributional split criteria, such as maximum mean discrepancy, in order to
adapt forest weights to the full conditional law of a response.  In univariate
tabular regression, however, the extra sensitivity is not automatically useful:
most finite samples appear to be dominated by location signal, while
higher-order distributional structure is weak or too noisy to estimate at split
time.  Across synthetic diagnostics, oracle comparisons, adaptive-frequency
variants, multivariate probes, and real-data benchmarks, we find that
CART-style splitting is the stronger default for univariate distributional
prediction.  The evidence supports a simple characterization: distributional
splitting helps only when non-location structure is both present and estimable;
otherwise it dilutes split-selection power away from the mean.
\end{abstract}

\section{Question}

Random forests for distributional prediction estimate not just a conditional
mean, but a conditional law.  Quantile regression forests do this with standard
CART tree structure and distributional leaf weights \citep{meinshausen2006}.
Distributional random forests instead propose splitting rules that compare the
full response distributions in candidate children \citep{cevid2022}.  This is a
natural generalization, especially for multivariate responses where a scalar
ordering or separate marginal quantile fits are unsatisfactory.

The narrower question studied here is deliberately less ambitious:
\emph{for univariate tabular regression, is distributional splitting a better
default than CART-style mean splitting?}  This note records a negative answer
with a positive interpretation.  CART-style splitting is difficult to beat in
this regime, not because distributional criteria are incoherent, but because the
finite-sample split-selection problem is usually mean dominated.

\section{Split Criteria as Signal Allocation}

For a node with observations $(X_i,Y_i)$, a candidate split $A \mid A^c$ is
chosen by maximizing an empirical separation score between the child responses.
CART uses the between-group mean separation
\[
  \frac{n_L n_R}{n^2}\|\bar Y_L-\bar Y_R\|^2,
\]
which is the identity-kernel mean-embedding score.  Random Fourier feature MMD
splitting replaces $Y$ by a complex feature map $\psi(Y)$ and scores
\[
  \frac{n_L n_R}{n^2}
  \left\| \frac{1}{n_L}\sum_{i\in L}\psi(Y_i)
        - \frac{1}{n_R}\sum_{i\in R}\psi(Y_i) \right\|^2 .
\]
With enough features and enough data this score can detect changes beyond the
mean.  At finite sample sizes it also spends split-selection capacity on noisy
higher-order coordinates.

This suggests the working characterization tested below.  CART concentrates on
the dominant location signal.  MMD spreads sensitivity across location and
higher-order distributional structure.  MMD should win only when the additional
structure is strong enough, or the sample size large enough, for its split-time
signal-to-noise ratio to exceed the cost of estimating it.

\section{Evidence}

The experiments were run as paired-seed comparisons: the train/test split,
forest seeds, tree hyperparameters, honesty setting, and cutpoint grid were held
fixed while varying only the split criterion.  Metrics are RMSE for the
conditional mean and CRPS / energy score for distributional prediction.

\paragraph{Synthetic diagnostics.}
The first set of experiments swept the univariate synthetic mechanisms used in
the DRF literature, plus targeted variants designed to isolate mean shifts,
scale shifts, quantile changes, and weak higher-moment effects.  CART dominated
or tied MMD in the regimes where the conditional law was primarily location
driven.  MMD became plausible only when higher-order structure was deliberately
made large enough to be visible at split time.

\begin{table}[h]
  \centering
  \caption{Synthetic univariate summary.  Replace this placeholder with paired
  mean differences, e.g. MMD minus CART, so negative CRPS/energy indicates an
  MMD improvement.}
  \begin{tabular}{lrrrr}
    \toprule
    dataset & honesty & $\Delta$RMSE & $\Delta$CRPS & $\Delta$energy \\
    \midrule
    paper\_quantile\_1 & 0.5 & TODO & TODO & TODO \\
    paper\_quantile\_2 & 0.5 & TODO & TODO & TODO \\
    paper\_quantile\_3 & 0.5 & TODO & TODO & TODO \\
    \bottomrule
  \end{tabular}
\end{table}

\paragraph{Adaptive and oracle checks.}
Adaptive-frequency MMD was used as a rescue attempt: if uniform RFF averaging
was hiding useful distributional signal, selecting the strongest coordinates at
each split should expose it.  Oracle and adaptive runs did not change the main
conclusion.  This points away from a simple tuning failure and toward the
finite-sample allocation explanation.

\paragraph{Real data.}
On univariate tabular benchmarks, CART-style splitting was again the stronger
default.  The exception worth isolating is California housing, where MMD can
look competitive.  That case confounds two explanations: the dataset is large,
and housing prices may contain estimable heteroskedastic structure.  A
subsample experiment over $n\in\{2000,4000,8000,16000\}$ separates these
readings.  If MMD improves only as $n$ grows, the evidence favors the
estimability threshold story; if it improves already at small $n$, the evidence
favors unusually strong distributional structure.

\begin{table}[h]
  \centering
  \caption{Real-data paired-seed summary.  Fill with means and standard errors
  for CART and MMD, or paired MMD--CART differences.}
  \begin{tabular}{lrrrrr}
    \toprule
    dataset & $n_{\mathrm{train}}$ & honesty & criterion & CRPS & fit time \\
    \midrule
    diabetes & TODO & 0.5 & CART & TODO & TODO \\
    diabetes & TODO & 0.5 & MMD & TODO & TODO \\
    california & TODO & 0.5 & CART & TODO & TODO \\
    california & TODO & 0.5 & MMD & TODO & TODO \\
    \bottomrule
  \end{tabular}
\end{table}

\section{Characterization}

The empirical pattern can be summarized as follows.

\begin{quote}
For univariate distributional forests, distributional splitting helps when
non-location conditional structure is present and estimable at node scale.  In
mean-dominated finite samples, CART-style splitting is the better default
because it allocates split-selection power to the strongest signal.
\end{quote}

This statement is intentionally scoped.  It is not a claim that distributional
forests are inferior to quantile regression forests, nor that MMD is a poor
criterion in general.  The core motivation for DRF is multivariate response
distributions, where CART/QRF-style univariate orderings are not a complete
solution.  The claim here is about the common univariate tabular setting: when
the target is scalar and the conditional law is mostly a shifted version of
itself, a distributional split criterion is often paying variance to chase weak
structure.

\section{Implications}

For practice, the recommendation is simple.  Use CART-style splitting as the
univariate default, then use distributional criteria as targeted tools for data
sets with clear heteroskedasticity, tail behavior, or non-location structure.
When distributional splitting is used, paired-seed comparisons and honesty
settings should be reported: in these experiments, honesty can be as important
as the criterion itself.

For research, the next useful theory is not a generic dominance theorem.  The
useful object is a finite-sample comparison of split-score signal-to-noise:
mean separation versus kernel mean separation under local alternatives.  Such a
calculation would turn the empirical characterization above into a sharp
threshold statement.

\section{Limitations}

The experiments here are mostly univariate.  They do not exhaust the design
space where DRF is most naturally motivated: multivariate responses, dependence
structure, and conditional distributional changes with little or no mean shift.
They also compare CART-style splitting to MMD splitting inside the same forest
implementation; a full external QRF package comparison is a separate baseline.

\bibliographystyle{plainnat}
\bibliography{main}

\end{document}
